\section{Methodology}
\label{sec:methodology}

\subsection{Workload, scenarios, and tooling}

We use \texttt{k6 v1.0}~\cite{k6} as the load generator. It is run on
a dedicated \texttt{c5.large} EC2 instance in
\texttt{us-east-1a}, the same availability zone as the targets, so that
network jitter to the ALB is bounded by intra-AZ RTT (sub-millisecond
in our pings). Each run begins with a thirty-second warm-up phase
during which we issue 50~serial requests against \texttt{/healthz} on
each platform to make sure the connection pools are full and any JIT
state is hot. Per-request timings are recorded in CSV and a JSON
summary (with bootstrapped percentile statistics) is recorded at the
end. The four scenarios are:

\begin{description}[leftmargin=1em,labelindent=0pt,itemsep=2pt]
\item[Steady-state RPS sweep.] Constant arrival rate at
      $\{50, 100, 250, 500\}$ RPS for 60\,s each.  We measure
      end-to-end latency for the GET-redirect path.
\item[Step-arrival burst.] One minute at 5\,RPS, three minutes at
      400\,RPS, one minute back at 5\,RPS. The transition is what we
      call ``elasticity'' and we estimate the time-to-steady-state
      from the trace (see below).
\item[Cold-start probe.] One virtual user issues 15 sequential GETs
      with 45\,s gaps, designed to force a fresh container init for
      Lambda. EC2 and Fargate are run with the same script as a
      sanity baseline.
\item[Mixed read/write.] 70\,\% GET-redirect, 30\,\% POST-shorten at
      300\,RPS for two minutes; representative of a real shortener.
\end{description}

\subsection{Statistics}

Per-request raw latency is highly non-normal (heavy right tail), so we
report distribution percentiles rather than mean$\pm$sd. Confidence
intervals on each percentile are computed by non-parametric bootstrap
with $B{=}1000$ resamples and $\alpha{=}0.05$. The error bars in
\cref{fig:steady} are these intervals.

\paragraph{Time-to-steady-state.} We define the elasticity of a
deployment as the elapsed time between the start of the burst and the
first one-second bin whose p95 latency falls within $1.2\times$ of the
\emph{post-transient steady value}, where the steady value is the
median of the last 20\,\% of the burst window. Whenever no such bin
exists --- which only happens for Lambda when burst RPS exceeds the
account concurrency limit --- we record $\infty$.

\paragraph{Cost model.} The cost analysis combines the AWS billing
formulas of \cref{sec:background} with our measured median execution
time. Concretely, for sustained $\rho$ RPS, our cost-per-month
(USD) is
\begin{equation}
\label{eq:cost}
\begin{aligned}
C_{\text{EC2}}      &= n(\rho)\cdot p_{\text{ec2}}\cdot T \\
C_{\text{Fargate}}  &= n(\rho)\cdot (v\cdot p_{\text{vCPU}} + g\cdot p_{\text{GB}})\cdot T \\
C_{\text{Lambda}}   &= \rho\cdot T\cdot
                       \bigl(p_{\text{req}} + g\cdot \overline{t}\cdot p_{\text{GB-s}}\bigr) + p_{\text{PC}}
\end{aligned}
\end{equation}

with $n(\rho)$ the steady-state number of compute units required to
sustain $\rho$ RPS at our measured saturation throughput, $T$ the
month-hours, $\overline{t}$ the median execution time of a request,
$g{=}1$\,GB and $v{=}0.5$ vCPU. Because DynamoDB and the ALB
contribute the same cost to all three paradigms, we report the
\emph{compute-only} cost when comparing them.
